\chapter{Слабые топологии и компактность.}
\section{Теорема Банаха-Алаоглу.}\label{seq:banach-alaogly}
\begin{definition}
    Пусть $(X, \tau)$ -- топологическое векторное пространство,
    Пусть $0 \in V \in \tau$.
    Определим $\Gamma(V)$ (так называемую поляру окрестности $V$) как
    \[
        \Gamma(V) = \{f \in (X, \tau)^* \mid |f(x)| \leq 1 \ \forall x \in V\},
    \]
    где $(X, \tau)^*$ понимается как топологическое сопряжённое, то есть множество всех $\tau$-непрерывных на $X$ функционалов.
\end{definition}
\begin{note}
    Заметим, что на топологическом сопряжённом $(X, \tau)^*$ мы можем ввести *-слабую топологию, так как $X$ разделяет точки $(X, \tau)^*$ просто по определению.
    По замечанию с первой лекции эта *-слабая топология $\tau_{w*}$ совпадает с топологией, индуцированной с $\Cm^X$ с топологией Тихонова на подпространство $\tau$-непрерывных функционалов.
\end{note}
\begin{example}
    Пусть $X$ -- нормированное пространство.
    Тогда $\Gamma(B_R(0)) = B_{1/R}^*[0]$, так как
    \[
        \begin{aligned}
            \Gamma(B_R(0)) &= \\
            &= \{f \in X^* \mid |f(x)| \leq 1 \ \forall \|x\| < R\} \\
            &= \left\{f \mid \left|f(x)\right| \leq \frac{1}{R} \  \forall \|x\| \leq 1\right\} \\
            &= \left\{f \mid \|f\| \leq \frac{1}{R}\right\} = \\
            &= B_{1/R}^*[0].
        \end{aligned}
    \]
\end{example}
\begin{theorem}[Банах, Алаоглу]
    Пусть $(X, \tau)$ -- топологическое векторное пространство, и $V$ -- $\tau$-окрестность нуля.
    Тогда $\Gamma(V)$ является *-слабым компактом в $((X, \tau)^*, \tau_{w*})$.
\end{theorem}
\begin{proof}
    Рассмотрим $(X, \tau)^* \subset \Cm^{X}$.
    В $\Cm^X$ введём топологию Тихонова $\tau_T$.
    Тогда $\tau_{w*}$ -- сужение $\tau_T$ из $\Cm^X$ на $(X, \tau)^*$.

    Мы хотим показать, что $\Gamma(V)$ вложено в некоторый компакт в $\Cm^X$ с $\tau_T$.
    Тогда, если мы покажем замкнутость $\Gamma(V)$ в $\tau_T$, то оно будет *-слабым компактом как замкнутое подмножество компакта.

    Для всякого $x \in X$ верно
    \[
        \Cm \ni 0 \cdot x = 0 \in V \subset X.
    \]
    Тогда существует $\delta_x > 0$ такое, что $\delta_x \cdot x \in V$ (в силу непрерывности умножения на скаляр в топологическом векторном пространстве).
    И, следовательно,
    \[
        |f(\delta_x x)| \leq 1 \ \forall f \in \Gamma(V).
    \]
    Тогда
    \[
        |f(x)| \leq \frac{1}{\delta_x}.
    \]
    Рассмотрим $K_x = \left\{\lambda \in \Cm \mid |\lambda| \leq \frac{1}{\delta_x}\right\}$.
    Тогда множество $M = \prod_{x \in X} K_x$ таково, что $\Gamma(V) \subset M$, а сам $M$ компактен в топологии Тихонова по теореме Тихонова.

    Покажем $\tau_T$-замкнутость $\Gamma(V)$ в $(\Cm^X, \tau_T)$.
    Пусть $h \in [\Gamma(V)]_{\tau_T} \subset \Cm^{X}$.
    Мы хотим показать, что $h \in (X, \tau)^*$, то есть $h\colon X \to \Cm$ является линейным и $\tau$-непрерывным, и $|h(x)| \leq 1$ при всяком $x \in V$.

    То, что $h \in [\Gamma(V)]_{\tau_T}$ означает, что $\forall U(h) \in \tau_T$ существует $f \in U(h) \cap \Gamma(V)$,

    Покажем линейность $h$.
    Пусть $x, y \in X$ и $\alpha, \beta \in \Cm$.
    Пусть $\epsilon > 0$.
    Рассмотрим $U(h) = V(h, x, \epsilon) \cap V(h, y, \epsilon) \cap V(h, \alpha x + \beta y, \epsilon)$,
    где
    \[
        V(h, z, \epsilon) = \{g \in \Cm^X \mid |g(z) - h(z)| < \epsilon\}.
    \]
    Существует $f \in U(h) \cap \Gamma(V)$
    Тогда
    \[
        \begin{cases}
            |h(x) - f(x)| < \epsilon, \\
            |h(y) - f(y)| < \epsilon, \\
            |h(\alpha x + \beta y) - f(\alpha x + \beta y)| < \epsilon.
        \end{cases}
    \]
    А значит
    \[
        |h(\alpha x + \beta y) - \alpha h(x) - \beta h(y)| \leq \epsilon(1 + |\alpha| + |\beta|).
    \]
    Так как это верно для произвольного $\epsilon$, то $h$ -- линеен.

    Покажем непрерывность $h$.
    Достаточно показать непрерывность $h$ в $0$.
    Заметим, что если показать, что $\forall \epsilon > 0$ существует $W_\epsilon(0) \in \tau$ такой, что $|h(x)| \leq \epsilon$ при всяком $x \in W_{\epsilon}(0)$, то $h$ будет непрерывным в $0$.
    Рассмотрим $W_\epsilon(0) = \epsilon V \in \tau$.
    Если верна оценка, что $|h(x)| \leq 1$ при $x \in V$, то такая $W_\epsilon(0)$ будет искомой и непрерывность будет показано.

    Покажем теперь эту оценку.
    Для всякого $x \in V$ и $\epsilon > 0$ рассмотрим окрестность $V(h, x, \epsilon)$.
    Тогда существует $f \in V(h, x, \epsilon) \cap \Gamma(V)$.
    А значит
    \[
        |h(x)| \leq |h(x) - f(x)| + |f(x)| < 1 + \epsilon \rightarrow 1, \epsilon \rightarrow +0.
    \]
    И таким образом показана нужная оценка, следовательно -- непрерывность функционала, а значит и замкнутость $\Gamma(V)$ в топологии Тихонова, что и даёт нам искомую *-слабую компактность.
\end{proof}
\begin{note}
    Заметим, что $B^*_r[0]$ в $(X, \|\cdot\|)^*$ не является *-слабым секвенциальным компактом.

    Пусть $X = \ell^{\infty}$.
    Рассмотрим $B_1^*[0] \subset (\ell^{\infty})^*$.
    В нём точно лежат координатные функционалы $f_n$ вида
    \[
        f_n(x) = x(n).
    \]
    Пусть у последовательности координатных функционалов существует *-слабо сходящаяся подпоследовательность $f_{n_k}$.
    Определим $x \in \ell^{\infty}$ как
    \[
        x(n) = \begin{cases}
        (-1)^k, & \exists k\colon n = n_k, \\
                   0, & \not\exists k\colon n = n_k.
        \end{cases}
    \]
    Тогда
    \[
        f_{n_k}(x) = (-1)^k,
    \]
    а значит $f_{n_k}$ не может быть *-слабо сходящейся последовательностью.

    Но если $X$ -- сепарабельно, то это верно (в силу метризуемости *-слабой топологии на шаре и критерия Фреше).
\end{note}
\begin{proposition}
    Пусть $X$ -- нормированное пространство и $\Phi$ -- каноническая изометрия Банаха.
    Пусть $L = \Phi(X)$.
    Тогда $\Phi$ является гомеоморфизмом между $(X, \tau_w)$ и $(L, \tau_{w^*})$, где $\tau_{w^*}$ -- сужение *-слабой топологии на $X^{**}$ на $L$.
\end{proposition}
\begin{proof}
    Пусть $G \in \tau_{w*}$ и $x \in \Phi^{-1}(G)$, то есть $\Phi_x \in G$.
    Тогда
    \[
        \Phi_x \in \bigcap_{k = 1}^{N} V_*(\Phi_x, f_k, \epsilon) \subset G,
    \]
    следовательно
    \[
        \tau_w \ni U(x) = \bigcap_{k = 1}^{N} V(x, f_k, \epsilon),
    \]
    поскольку \[y \in U(x) \Longleftrightarrow |f_k(x) - f_k(y)| <\epsilon \Longleftrightarrow |\Phi_x(f_k) - \Phi_y(f_k)| < \epsilon \Longleftrightarrow \Phi_y \in \bigcap_{k = 1}^{N} V_*(\Phi_x, f_k, \epsilon).\]
    То есть $U(x) \subset \Phi^{-1}(G)$ при всяком $x$, и, следовательно $\Phi^{-1}(G) \in \tau_w$ что и завершает доказательство того, что $\Phi$ -- гомеоморфизм (то, что $\Phi(G)$ открыто в $L$ доказывается абсолютно аналогично).
\end{proof}
\begin{corollary}
    Пусть $X$ -- рефлексивно и $\Phi$ -- каноническая изометрия Банаха.
    Тогда, по теореме Банаха-Алаоглу, $B_1^{**}[0]$ -- компакт в *-слабой топологии на $X^{**}$.
    Заметим, что слабая топология на $X$ и *-слабая топология на $X^{**}$ порождается одним и тем же набором функционалов, а значит они совпадают, а $\Phi$ является гомеоморфизмом между ними.
    Тогда $B_1[0]$ является слабо компактным в $X$.
\end{corollary}
\begin{note}
    Если $X$ -- рефлексивно и сепарабельно, то $X^{**} = \Phi(X)$ -- тоже сепарабельно, а значит и $X^*$ -- сепарабельно, и тогда слабая топология в $X$ является метрической на шаре в $X$, а значит он является и слабым секвенциальным компактом.
\end{note}
\begin{note}[Семинарский факт]
    Заметим, что если $X$ -- нормированное пространство, и $\Phi$ -- каноническая изометрия Банаха, то
    \[
        [\Phi(X)]_{\tau_{w*}} = X^{**}.
    \]
\end{note}
\begin{note}[Следствие из задачи из задания]
    Если $X$ -- рефлексивное пространство, то и $X^*$ -- рефлексивно.
    А значит тогда в $X^*$ слабая и *-слабая топологии совпадают, и шар в сопряжённом пространстве будет слабым компактом.
\end{note}

\begin{theorem}\label{th:hilbert-reflexivity}
    Пусть $H$ -- гильбертово пространство.
    Тогда $H$ -- рефлексивно.
\end{theorem}
\begin{proof}
    Пусть $F \in H^{**}$.
    Необходимо показать сюръективность канонического вложения Банаха, то есть то, что существует $x \in H$ такой, что
    \[
        F(f) = f(x) \quad \forall f \in H^*.
    \]
    В силу теоремы Рисса $\forall f \in H^*$ существует $y_f \in H\colon f(x) = \langle x, y_f\rangle$. Изометрию Рисса обозначим как $R\colon H \to H^*$.
    Рассмотрим отображение $\phi_F\colon H \to \Cm$, определённое как
    \[
        \phi_F(y) = F(Ry).
    \]
    Так как $F$ -- линейно, $R$ -- антилинейно, то $\phi_F$ -- антилинейно и, следовательно, $\overline{\phi_F}$ -- линейный функционал.
    Опять воспользуемся теоремой Рисса и получим $y_\phi\colon \overline{\phi_F}(z) = \langle z, y_\phi \rangle$, то есть $\phi_F(z) = \langle y_\phi, z \rangle$.
    Рассмотрим произвольный $f \in H^*\colon f = Rh$.
    Тогда
    \[
        F(f) = F(Rh) = \phi_F(h) = \langle y_\phi, h \rangle = f(y_\phi),
    \]
    и $F = J_H y_\phi$ -- что и даёт сюръективность изометрии Банаха, а значит и рефлексивность $H$.
\end{proof}